\documentclass{article}

\usepackage{hyperref}
\usepackage{graphicx}
\usepackage{fancyhdr}
\usepackage{extramarks}
\usepackage{amsmath}
\usepackage{cleveref}
\usepackage{amsthm}
\usepackage{amsfonts}
\usepackage{tikz}
\usepackage[plain]{algorithm}
\usepackage{algpseudocode}

\usetikzlibrary{automata,positioning}

%
% Basic Document Settings
%

\topmargin=-0.45in
\evensidemargin=0in
\oddsidemargin=0in
\textwidth=6.5in
\textheight=9.0in
\headsep=0.25in

\linespread{1.1}

\pagestyle{fancy}
\lhead{\hmwkAuthorName}
\chead{\hmwkClass , \hmwkClassInstructor: \hmwkTitle}
\rhead{\firstxmark}
\lfoot{\lastxmark}
\cfoot{\thepage}

\renewcommand\headrulewidth{0.4pt}
\renewcommand\footrulewidth{0.4pt}

\setlength\parindent{2em}
\setlength{\parskip}{0.5em}

%
% Create Problem Sections
%

\newcommand{\enterProblemHeader}[1]{
    \nobreak\extramarks{}{Problem \arabic{#1} continued on next page\ldots}\nobreak{}
    \nobreak\extramarks{Problem \arabic{#1} (continued)}{Problem \arabic{#1} continued on next page\ldots}\nobreak{}
}

\newcommand{\exitProblemHeader}[1]{
    \nobreak\extramarks{Problem \arabic{#1} (continued)}{Problem \arabic{#1} continued on next page\ldots}\nobreak{}
    \stepcounter{#1}
    \nobreak\extramarks{Problem \arabic{#1}}{}\nobreak{}
}

\setcounter{secnumdepth}{0}
\newcounter{partCounter}
\newcounter{homeworkProblemCounter}
\setcounter{homeworkProblemCounter}{1}
\nobreak\extramarks{Problem \arabic{homeworkProblemCounter}}{}\nobreak{}

%
% Homework Problem Environment
%
% This environment takes an optional argument. When given, it will adjust the
% problem counter. This is useful for when the problems given for your
% assignment aren't sequential. See the last 3 problems of this template for an
% example.
%
\newenvironment{homeworkProblem}[1][-1]{
    \ifnum#1>0
        \setcounter{homeworkProblemCounter}{#1}
    \fi
    \section{Problem \arabic{homeworkProblemCounter}}
    \setcounter{partCounter}{1}
    \enterProblemHeader{homeworkProblemCounter}
}{
    \exitProblemHeader{homeworkProblemCounter}
}

%
% Homework Details
%   - Title
%   - Due date
%   - Class
%   - Section/Time
%   - Instructor
%   - Author
%

\newcommand{\hmwkTitle}{Problem Set 3}
\newcommand{\hmwkClass}{CS229}
\newcommand{\hmwkClassInstructor}{Prof. Andrew Ng}
\newcommand{\hmwkAuthorName}{\textbf{Connor Yan}}

%
% Title Page
%

\title{
    \vspace{2in}
    \textmd{\textbf{\hmwkClass}}\\
    \textmd{\textbf{\hmwkTitle}}\\
    \vspace{3in}
}

\author{\hmwkAuthorName}
\date{\today}

\renewcommand{\part}[1]{\textbf{\large Part \Alph{partCounter}}\stepcounter{partCounter}\\}

% Alias for the Solution section header
\newcommand{\solution}{\textbf{\large Solution}}

% Probability commands: Expectation, Variance, Covariance, Bias
\newcommand{\E}{\mathbb{E}}
\newcommand{\Var}{\mathrm{Var}}
\newcommand{\Cov}{\mathrm{Cov}}
\newcommand{\Bias}{\mathrm{Bias}}
\newcommand{\dd}{\mathrm{d}}
\newcommand{\R}{\mathbb{R}}
\newcommand{\mb}{\mathrm}

\begin{document}

\maketitle

\pagebreak

\begin{homeworkProblem}
\begin{enumerate}
    \item[(a)] The network implies
    \[
        h_2^{(i)} = g\left( (w^{[1]}_{,2})^T x^{(i)} + w^{[1]}_{0,2} \right), \ 
        o^{(i)} =  g\left((w^{[2]})^T h^{(i)} + w^{[2]}_{0} \right), \ \text{where } g(z) = \frac{1}{1 + \exp(-z)}
    \]
    The derivative of loss function w.r.t. $w^{[1]}_{1,2}$ is
    \[
    \begin{aligned}
        \frac{\partial \ell}{\partial w^{[1]}_{1,2}}
        &= \frac{\partial \ell}{\partial o} \cdot \frac{\partial o}{\partial h_2} \cdot \frac{\partial h_2}{\partial w^{[1]}_{1,2}} \\
        &= \frac{2}{m} \sum_{i=1}^m (o^{(i)}-y^{(i)}) \cdot o^{(i)}(1-o^{(i)})w^{[2]}_{2} \cdot h_2^{(i)}(1-h_2^{(i)})x^{(i)}
    \end{aligned}
    \]
    Hence, the GD update is
    \[
    \begin{aligned}
        w^{[1]}_{1,2} &:= w^{[1]}_{1,2} - \alpha \frac{\partial \ell}{\partial w^{[1]}_{1,2}} \\
        &= w^{[1]}_{1,2} - \alpha \cdot \frac{2}{m}\sum_{i=1}^m(o^{(i)}-y^{(i)}) o^{(i)}(1-o^{(i)})w^{[2]}_{2} h_2^{(i)}(1-h_2^{(i)})x^{(i)}
    \end{aligned}
    \]
    
    \item[(b)] The two classes are seperable by a triangle. The constraint indicating that a data point lies inside of one side of the triangle can be written as
    \[
        w_{1,j}^{[1]} x^{(i)}_1 + w_{2,j}^{[1]}x_2^{(i)} + w^{[1]}_{0,j} \ge 0
    \]
    That is to say, if weights and bias are rightly set, given a data point is inside of the triangle,
    \[
        h_j = f(w_{1,j}^{[1]} x^{(i)}_1 + w_{2,j}^{[1]}x_2^{(i)} + w^{[1]}_{0,j}) = 1
    \]
    Hence, to achieve 100\% accuracy is to ensure $o=0$ only if all 3 neurons in the hidden layer to output 1 namely, the output neuron is something like
    \[
        o = f(2.5 - h_1 - h_2 - h_3),
    \]
    so that $o=0$ if and only if the data point is inside of all three sides of the triangle.

    A set of weights that can ensure 100\% accuracy is
    \[
    \begin{aligned}
        w^{[1]}_{0,1}&=-0.5, & w^{[1]}_{1,1}&=1, & w^{[1]}_{2,1}&=0,\\
        w^{[1]}_{0,2}&=-0.5, & w^{[1]}_{1,2}&=0, & w^{[1]}_{2,2}&=1,\\
        w^{[1]}_{0,3}&=4,    & w^{[1]}_{1,3}&=-1,& w^{[1]}_{2,3}&=-1,\\[0.5em]
        w^{[2]}_{0}&=2.5, & w^{[2]}_{1}&=-1, &
        w^{[2]}_{2}&=-1, & w^{[2]}_{3}&=-1.
    \end{aligned}
    \]

    \item[(c)] Impossible to have 100\% accuracy. With the linear activation functions $g(x) = x$ for $h_1$, $h_2$, and $h_3$, and step function $f(x)$ for $o$, the output is simply
    \[
    \begin{aligned}
        o &= f(w^{[2]}_0 + \sum_{j=1}^3w^{[2]}_jh_j) \\
        &= f(w^{[2]}_0 + \sum_{j=1}^3w^{[2]}_j(w^{[1]}_{0,j} + \sum_{i=1}^2 w^{[1]}_{i,j}x_i)) \\
        &= f(w_0 + w_1x_1 + w_2x_2)
    \end{aligned}
    \]
    The neural network reduces to a linear classifier. However, the given dataset is not linear-separable, so it is impossible to have 100\% accuracy.
    
\end{enumerate}
\end{homeworkProblem}

\begin{homeworkProblem}
\begin{enumerate}
    \item[(a)]
    \[
    \begin{aligned}
        D_\mathrm{KL}(P\|Q) &= \sum_{x\in\mathcal{X}} P(x) \log \frac{P(x)}{Q(x)}\\
        &= \E_{X\sim P}\left[ \log \frac{P(X)}{Q(X)} \right] \\
        &= \E_{X\sim P}\left[ -\log \frac{Q(X)}{P(X)} \right] \\
    \end{aligned}
    \]
    Since $f(x) = -\log x$ is convex, according to Jensen's inequality,
    \[
    \begin{aligned}
        D_\mathrm{KL}(P||Q) = \E_{X\sim P}\left[ -\log \frac{Q(X)}{P(X)} \right]
        &\ge -\log \E_{X\sim P}\left[ \frac{Q(X)}{P(X)} \right] \\
        &= -\log \sum_{x\in\mathcal{X}} Q(x) \\
        &= -\log 1 \\
        &= 0
    \end{aligned}
    \]
    Note that the inequality has $\E_{X\sim P}\left[ -\log \frac{Q(X)}{P(X)} \right] = -\log \E_{X\sim P}\left[ \frac{Q(X)}{P(X)} \right]$ if and only if 
    \[
    \begin{aligned}
        \frac{Q(X)}{P(X)} = \E_{X\sim P}\left[ \frac{Q(X)}{P(X)} \right] = 1,
    \end{aligned}
    \]
    i.e. $Q=P$.
    Hence,
    \[
    \forall P,Q \ \ D_\mathrm{KL}(P\|Q)\ge 0 \text{ and } D_\mathrm{KL}(P\|Q) = 0 \text{ iif.} P=Q
    \]

    \item[(b)]
    \[
    \begin{aligned}
        D_\mathrm{KL}(P(X,Y)\|Q(X,Y))
        &= \sum_x\sum_y P(x,y) \log \frac{P(x,y)}{Q(x,y)} \\
        &= \sum_x\sum_y P(y|x)P(x) \log \frac{P(y|x)P(x)}{Q(y|x)Q(x)} \\
        &= \sum_x P(x) \sum_y P(y|x) \left( \log\frac{P(y|x)}{Q(y|x)} + \log\frac{P(x)}{Q(x)} \right) \\
        &= \sum_x P(x) \sum_y P(y|x) \log\frac{P(y|x)}{Q(y|x)} 
           + \sum_x P(x) \log\frac{P(x)}{Q(x)} \sum_y P(y|x) \\
        &= D_\mathrm{KL}(P(Y|X)\|Q(Y|X)) + D_\mathrm{KL}(P(X)\|Q(X))
    \end{aligned}
    \]

    \item[(c)]
    \[
    \begin{aligned}
        D_\mathrm{KL}(\hat{P}\|P_\theta)
        &= \sum_{i=1}^m \hat{P}(x^{(i)}) \log{\frac{\hat{P}(x^{(i)})}{P_\theta(x^{(i)})}} \\
        &= H(\hat{P}) - \sum_{i=1}^m \hat{P}(x^{(i)}) \log{P_\theta(x^{(i)})} \\
        &= H(\hat{P}) - \sum_{i=1}^m \log{P_\theta(x^{(i)})} \cdot \frac{1}{m}\sum_{j=1}^m{ 1\{x^{(j)} = x^{(i)}\} } \\
        &= H(\hat{P}) - \frac{1}{m} \sum_{i=1}^m \log{P_\theta(x^{(i)})}
    \end{aligned}
    \]
    Hence,
    \[
    \arg\min_\theta D_\mathrm{KL}(\hat{P}\|P_\theta) = \arg\max_\theta \sum_{i=1}^m \log{P_\theta(x^{(i)})}
    \]
    
\end{enumerate}
\end{homeworkProblem}

\begin{homeworkProblem}
\begin{enumerate}
    \item[(a)] 
    \[
    \begin{aligned}
        \E_{y\sim p(y;\theta)}[\nabla_{\theta'}\log{p(y;\theta')}|_{\theta'=\theta}]
        &= \int_{-\infty}^\infty{ p(y;\theta) \nabla_{\theta'}\log{p(y;\theta')}|_{\theta'=\theta} }  dy \\
        &= \int_{-\infty}^\infty{ \nabla_{\theta'} p(y;\theta')|_{\theta'=\theta} } dy \\
        &= \left. \nabla_{\theta'} \int_{-\infty}^\infty{ p(y;\theta') } dy \right|_{\theta'=\theta} \\
        &= \nabla_{\theta'} 1 |_{\theta'=\theta} \\
        &= 0
    \end{aligned}
    \]

    \item[(b)] Let
    \[
    \mu = \E_{y\sim p(y;\theta)}[\nabla_{\theta'}\log{p(y;\theta')}|_{\theta'=\theta}] = 0
    \]
    Then,
    \[
    \begin{aligned}
        \mathcal{I}(\theta) &= \Cov_{y\sim p(y;\theta)}[\nabla_{\theta'}\log{p(y;\theta')}|_{\theta'=\theta}] \\
        &= \E_{y\sim p(y;\theta)}[(\nabla_{\theta'}\log{p(y;\theta')}-\mu) (\nabla_{\theta'}\log{p(y;\theta')}-\mu)^T|_{\theta'=\theta}], \\
        &= \E_{y\sim p(y;\theta)}[\nabla_{\theta'}\log{p(y;\theta')} \nabla_{\theta'}\log{p(y;\theta')}^T|_{\theta'=\theta}]
    \end{aligned}
    \]

    \item[(c)] 
    \[
    \begin{aligned}
        \nabla_{\theta'}\log{p(y;\theta')} &= \frac{1}{p}\nabla_{\theta'}p(y;\theta')\\
        \nabla_{\theta'}^2\log{p(y;\theta')}
        &= -\frac{1}{p^2}\nabla_{\theta'}p(y;\theta')\left(\nabla_{\theta'}p(y;\theta')\right)^T + \frac{1}{p}\nabla_{\theta'}^2p(y;\theta')
    \end{aligned}
    \]
    \[
    \begin{aligned}
        \E_{y\sim p(y;\theta)}[-\nabla_{\theta'}^2\log{p(y;\theta')}|_{\theta'=\theta}]
        &= \int_{-\infty}^\infty{ -p\nabla_{\theta'}^2\log{p(y;\theta')}|_{\theta'=\theta} } dy \\
        &= \int_{-\infty}^\infty{ \left. \frac{1}{p}\nabla_{\theta'} p(y;\theta') \left(\nabla_{\theta'}p(y;\theta')\right)^T - \nabla_{\theta'}^2p(y;\theta') \right|_{\theta'=\theta} } dy \\
        &= \int_{-\infty}^\infty{ \left. p\nabla_{\theta'} \log p(y;\theta') \left(\nabla_{\theta'}\log p(y;\theta')\right)^T\right|_{\theta'=\theta} } dy + \nabla_{\theta'}^2 \int_{-\infty}^\infty{ p(y;\theta') } dy |_{\theta'=\theta} \\
        &= \E_{y\sim p(y;\theta)}[\nabla_{\theta'}\log{p(y;\theta')} \nabla_{\theta'}\log{p(y;\theta')}^T|_{\theta'=\theta}]\\
        &= \mathcal{I}(\theta)
    \end{aligned}
    \]

    \item[(d)] Let $\theta + d = \tilde{\theta}$, $f(\tilde{\theta}) = D_\mathrm{KL}(p_{\theta}\|p_{\tilde{\theta}})$.
    \[
    f(\theta) = D_\mathrm{KL}(p_{\theta}\|p_{\theta}) = 0
    \]
    \[
    \begin{aligned}
        \nabla_{\theta'}f(\theta') &= \nabla_{\theta'} D_\mathrm{KL}(p_{\theta}\|p_{\theta'}) \\
        &= \nabla_{\theta'} \int_{-\infty}^\infty{p_\theta(x) \log\frac{p_\theta(x)}{p_{\theta'}(x)}}dx \\
        &= -\nabla_{\theta'} \int_{-\infty}^\infty{p_\theta(x) \log p_{\theta'}(x)}dx \\
        &= -\E_{x\sim p(x)}[\nabla_{\theta'} \log p_{\theta'}(x)]
    \end{aligned}
    \]
    \[
    \begin{aligned}
        \nabla_{\theta'}^2 f(\theta') &= -\nabla_{\theta'} \int_{-\infty}^\infty{p_\theta(x) \nabla_{\theta'} \log p_{\theta'}(x)}dx \\
        &= \E_{x\sim p(x)}[-\nabla_{\theta'}^2 \log p_{\theta'}(x)]
    \end{aligned}
    \]
    The Taylor series of $f(\tilde{\theta})$ near $\theta$ is
    \[
    \begin{aligned}
        f(\tilde{\theta})
        &\approx f(\theta) + (\tilde{\theta}-\theta)^T\nabla_{\theta'}f(\theta')|_{\theta'=\theta} + \frac{1}{2}(\tilde{\theta}-\theta)^T(\nabla_{\theta'}^2f(\theta')|_{\theta'-\theta})(\tilde{\theta}-\theta) \\
        &=  0 - d^T\E_{x\sim p(x)}[\nabla_{\theta'} \log p_{\theta'}(x)|_{\theta'=\theta}] + \frac{1}{2}d^T\E_{x\sim p(x)}[-\nabla_{\theta'}^2 \log p_{\theta'}(x)|_{\theta'=\theta}]]d \\
        &= \frac{1}{2}d^T\mathcal{I}(\theta)d
    \end{aligned}
    \]

    \item[(e)] 
    \[
    \begin{aligned}
        \mathcal{L}(d,\lambda) &= \ell(\theta+d) - \lambda[D_\mathrm{KL}(p_\theta\|p_{\theta+d})-c] \\
        &\approx \ell(\theta) + d^T\nabla_{\theta'}\ell(\theta')|_{\theta'=\theta} - \lambda\left[ \frac{1}{2}d^T\mathcal{I}(\theta)d - c \right]
    \end{aligned}
    \]
    \[
    \nabla_d\mathcal{L}(d,\lambda) = \nabla_{\theta'}\ell(\theta')|_{\theta'=\theta} - \lambda \mathcal{I}^T(\theta)d = \nabla_{\theta'}\ell(\theta')|_{\theta'=\theta} - \lambda \mathcal{I}(\theta)d
    \]
    \[
    \nabla_\lambda\mathcal{L}(d,\lambda) = \frac{1}{2}d^T\mathcal{I}(\theta)d - c
    \]
    Hence, we have system of linear equations:
    \begin{align}
        \nabla_{\theta'}\ell(\theta')|_{\theta'=\theta} - \lambda \mathcal{I}^T(\theta)d &= 0 \tag{a}\\
        \frac{1}{2}d^T\mathcal{I}(\theta)d - c &= 0 \tag{b}
    \end{align}
    From (a) we have
    \[
    d = \frac{1}{\lambda} \mathcal{I}^{-1}(\theta) \nabla_{\theta'}\ell(\theta')|_{\theta'=\theta}
    \]
    Plug it into (b):
    \[
    \begin{aligned}
        \frac{1}{\lambda^2} (\nabla_{\theta'}\ell(\theta')|_{\theta'=\theta})^T \mathcal{I}^{-1}(\theta) \mathcal{I}(\theta) \mathcal{I}^{-T}(\theta) \nabla_{\theta'}\ell(\theta')|_{\theta'=\theta} &= 2c \\
        \frac{1}{\lambda^2} (\nabla_{\theta'}\ell(\theta')|_{\theta'=\theta})^T \mathcal{I}^{-T}(\theta) \nabla_{\theta'}\ell(\theta')|_{\theta'=\theta} &= 2c \\
        \lambda &= \sqrt{\frac{1}{2c} (\nabla_{\theta'}\ell(\theta')|_{\theta'=\theta})^T \mathcal{I}^{-1}(\theta) \nabla_{\theta'}\ell(\theta')|_{\theta'=\theta}}
    \end{aligned}
    \]
    Plug it back into (a), which yields natural gradient:
    \[
    \boxed{
        d^* = \sqrt{\frac{2c}{(\nabla_{\theta'}\ell(\theta')|_{\theta'=\theta})^T \mathcal{I}^{-1}(\theta) \nabla_{\theta'}\ell(\theta')|_{\theta'=\theta}} } \mathcal{I}^{-1}(\theta) \nabla_{\theta'}\ell(\theta')|_{\theta'=\theta}
    }
    \]

    \item[(f)]
    \[
    \begin{aligned}
        \mathcal{I}(\theta)
        &= \E_{y\sim p(y;\theta)}[-\nabla_{\theta'}^2\log{p(y;\theta')}|_{\theta'=\theta})] \\
        &= \frac{1}{m} \sum_{i=1}^m \E_{y^{(i)}\sim p(y;\theta)}[-\nabla_{\theta'}^2\log{p(y^{(i)};\theta')}|_{\theta'=\theta}] \\
        &= \frac{1}{m} \E_{y^{(i)}\sim p(y;\theta)}[-\nabla_{\theta'}^2 \sum_{i=1}^m \log{p(y^{(i)};\theta')}|_{\theta'=\theta}] \\
        &= \frac{1}{m} \E_{y\sim p(y;\theta)}[-\nabla_{\theta'}^2 \ell(\theta') |_{\theta'=\theta}] \\
        &= \frac{1}{m} \E_{y\sim p(y;\theta)}[-\nabla_{\theta'}^2 \ell(\theta') \nabla_{\theta'}^2 \ell(\theta') |_{\theta'=\theta}] \nabla_{\theta'}^2 \ell(\theta') |_{\theta'=\theta} \\
        &= \frac{1}{m} \E_{y\sim p(y;\theta)}[-1] \nabla_{\theta'}^2 \ell(\theta') |_{\theta'=\theta} \\
        &= -\frac{1}{m} \nabla_{\theta'}^2 \ell(\theta') |_{\theta'=\theta}
    \end{aligned}
    \]
    The natural gradient is
    \[
    \begin{aligned}
        \tilde{d} &= \frac{1}{\lambda} \mathcal{I}^{-1}(\theta) \nabla_{\theta'}\ell(\theta')|_{\theta'=\theta} \\
        &= -\frac{1}{\lambda m} (\nabla_{\theta'}^2 \ell(\theta'))^{-1} \nabla_{\theta'}\ell(\theta')|_{\theta'=\theta}
    \end{aligned}
    \]
    Newton's method has update rule
    \[
    \begin{aligned}
        \theta &:= \theta - (\nabla^2_{\theta'}\ell(\theta')|_{\theta'=\theta})^{-1} \nabla_{\theta'}\ell(\theta')|_{\theta'=\theta} \\
        &= \theta + \lambda m \tilde{d}
    \end{aligned}
    \]
    
    
\end{enumerate}
\end{homeworkProblem}

\end{document}
